\documentclass[10pt]{beamer}

% Tema Limpo e Institucional IFCE Taua
\usetheme{Madrid}
\usecolortheme{whale}

% Pacotes Essenciais
\usepackage[utf8]{inputenc}
\usepackage[T1]{fontenc}
\usepackage[brazilian]{babel}
\usepackage{graphicx}
\usepackage{booktabs}
\usepackage{amsmath}
\usepackage{amssymb}
\usepackage{listings}
\usepackage{xcolor}
\usepackage{tikz}
\usetikzlibrary{shapes.geometric,arrows.meta,positioning,calc}

% Configuracao de Cores e Codigo (Python)
\definecolor{codegreen}{rgb}{0,0.6,0}
\definecolor{codegray}{rgb}{0.5,0.5,0.5}
\definecolor{codepurple}{rgb}{0.58,0,0.82}
\definecolor{backcolour}{rgb}{0.96,0.97,0.98}

\lstset{
    language=Python,
    backgroundcolor=\color{backcolour},
    basicstyle=\ttfamily\scriptsize,
    keywordstyle=\color{blue}\bfseries,
    stringstyle=\color{red!80!black},
    commentstyle=\color{green!60!black},
    numberstyle=\tiny\color{codegray},
    breaklines=true,
    showstringspaces=false,
    tabsize=4,
    frame=single,
    rulecolor=\color{blue!20}
}

% Metadados da Apresentacao
\title[Engenharia de Software -- Aula 24]{Engenharia de Software -- Aula 24}
\subtitle{Sprint 04: Casos de Uso, Diagrama de Classes e Início da Codificação}
\author[Prof. Me. Reginaldo Fernandes]{Prof. Me. Reginaldo Pereira Fernandes}
\institute[IFCE Campus Tauá]{
    Instituto Federal de Educação, Ciência e Tecnologia do Ceará -- Campus Tauá\\
    Técnico em Informática para Internet (TSII.210)
}
\date{Semana 1 -- Encontro 1 (100 min)}

\begin{document}

% Slide 1: Titulo
\begin{frame}
    \titlepage
\end{frame}

% Slide 2: Agenda do Encontro
\begin{frame}{Agenda do Encontro 1 (100 Minutos)}
    \tableofcontents
\end{frame}

\section{Contexto da Sprint e Liderança Técnica}

\begin{frame}{Cenário Corporativo: O Sistema BiblioTauá}
    \begin{itemize}
        \item \textbf{O Sistema:} Livraria Universitária e Biblioteca do IFCE Campus Tauá.
        \item \textbf{Status Atual (v0.3.0):}
        \begin{itemize}
            \item Módulo de Catálogo de Livros e Estoque funcional.
            \item Cadastro de Clientes e Discentes ativo.
            \item Persistência relacional em banco de dados SQLite.
        \end{itemize}
        \item \textbf{O Desafio da Sprint 04:}
        \begin{itemize}
            \item Feira do Livro do IFCE se aproximando.
            \item Demanda de negócio: conceder \textbf{Cupons Promocionais} aos estudantes.
            \item Meta: Entregar o Módulo de Cupons em 2 semanas (4 encontros).
        \end{itemize}
    \end{itemize}
\end{frame}

\begin{frame}{Sprint Planning: Fala da Tech Lead Aline Mendes}
    \begin{block}{Diretriz da Liderança Técnica}
    \textit{"Mariana assume a interface web. Lucas adapta o backend de pagamentos para o PIX. E \textbf{VOCÊ} assume a tarefa principal: \textbf{TASK-204 - Motor de Validação e Aplicação de Cupons}."}
    \end{block}
    \vspace{0.3cm}
    \begin{alertblock}{Ritmo Ágil de Alta Produtividade}
    Neste primeiro encontro, finalizaremos os diagramas de Casos de Uso e Classes e \textbf{já iniciaremos a codificação em Python}!
    \end{alertblock}
\end{frame}

\begin{frame}{O Ticket da Sprint: TASK-204}
    \begin{table}[h]
    \centering
    \footnotesize
    \begin{tabular}{ll}
        \toprule
        \textbf{Campo} & \textbf{Descrição do Ticket} \\
        \midrule
        Identificador & \textbf{TASK-204} \\
        Título & Motor de Validação e Aplicação de Cupons no Checkout \\
        Responsável & Desenvolvedor Júnior (Discente) \\
        Prioridade & Alta / Bloqueante para a Sprint 04 \\
        Estimativa & 4 Encontros de Trabalho (8 Horas-Aula) \\
        Branch Git & \texttt{feature/task-204-motor-cupons} \\
        Entrega Hoje & DUC $\rightarrow$ Diagrama de Classes $\rightarrow$ Exceções \& Modelo Cupom \\
        \bottomrule
    \end{tabular}
    \end{table}
\end{frame}

\section{Diagrama de Casos de Uso (UML 2)}

\begin{frame}{Modelagem Comportamental: Diagrama de Casos de Uso}
\centering
\resizebox{0.96\textwidth}{!}{%
\begin{tikzpicture}[>=Latex, font=\sffamily]
    % Fronteira do Sistema
    \draw[rounded corners=8pt, draw=blue!80!black, thick, fill=blue!3] (1.8,-3.3) rectangle (8.2,2.5);
    \node[anchor=north west, font=\scriptsize\bfseries, text=blue!80!black] at (2.0,2.3) {Fronteira: BiblioTauá Livraria Web};
    
    % Ator Stick Figure
    \node[draw, circle, minimum size=0.55cm, thick, fill=white] (head) at (0.4, 0.0) {};
    \draw[thick] (head.south) -- ++(0,-0.65) coordinate (torso);
    \draw[thick] ($(head.south)!0.3!(torso)$) -- ++(-0.4,-0.2);
    \draw[thick] ($(head.south)!0.3!(torso)$) -- ++(0.4,-0.2);
    \draw[thick] (torso) -- ++(-0.35,-0.65);
    \draw[thick] (torso) -- ++(0.35,-0.65);
    \node[below=0.75cm of torso, font=\scriptsize\bfseries, align=center] {Cliente /\\Estudante};

    % Casos de Uso (Perfeitamente distribuidos verticalmente sem sobreposicao)
    \node[draw=blue!80!black, fill=blue!10, ellipse, text width=2.8cm, align=center, font=\scriptsize] (uc01) at (5.0, 1.5) {UC01: Consultar\\Acervo de Livros};
    \node[draw=blue!80!black, fill=blue!10, ellipse, text width=2.8cm, align=center, font=\scriptsize] (uc02) at (5.0, 0.1) {UC02: Finalizar\\Pedido de Livros};
    \node[draw=green!60!black, fill=green!15, ellipse, text width=2.8cm, align=center, font=\scriptsize\bfseries] (uc04) at (5.0, -1.3) {UC04: Aplicar Cupom\\de Desconto};
    \node[draw=blue!80!black, fill=blue!10, ellipse, text width=2.8cm, align=center, font=\scriptsize] (uc05) at (5.0, -2.6) {UC05: Validar Regras\\de Elegibilidade};

    % Associacoes do Ator (Linhas diretas sem cruzamentos)
    \draw[thick] (head.east) -- (uc01.west);
    \draw[thick] (head.east) -- (uc02.west);
    \draw[thick] (head.east) -- (uc04.west);

    % Relacionamentos UML
    \draw[->, dashed, thick, draw=orange!90!black] (uc04) -- node[midway, right, font=\tiny\bfseries] {$\ll$extend$\gg$} (uc02);
    \draw[->, dashed, thick, draw=blue!80!black] (uc04) -- node[midway, right, font=\tiny\bfseries] {$\ll$include$\gg$} (uc05);

    % Caixa de Destaque / Guia de Leitura a Direita
    \node[draw=gray!50, fill=white, rounded corners=6pt, text width=3.8cm, align=left, inner sep=6pt, anchor=west] at (8.6, -0.4) {
        \textbf{\footnotesize Guia de Leitura UML:}\\[0.15cm]
        \scriptsize
        \textbf{\color{orange!90!black}$\ll$extend$\gg$ (Extensão):}\\
        Condicional: o cliente só aplica cupom se desejar durante o checkout.\\[0.2cm]
        \textbf{\color{blue!80!black}$\ll$include$\gg$ (Inclusão):}\\
        Obrigatório: ao aplicar cupom, o sistema valida compulsoriamente as regras.\\[0.2cm]
        \textbf{\color{green!60!black}Foco da TASK-204:}\\
        Motor de regras do \textbf{UC04}.
    };
\end{tikzpicture}
}
\end{frame}

\begin{frame}{Dissecando a DUC UC04: Regras de Negócio e Exceções}
    \begin{table}[h]
    \centering
    \tiny
    \begin{tabular}{lll}
        \toprule
        \textbf{Fluxo} & \textbf{Cenário / Regra de Negócio} & \textbf{Exceção Disparada em Python} \\
        \midrule
        \textbf{Principal} & Cupom ativo, dentro do prazo, cota disponível & Desconto calculado e total abatido \\
        \textbf{FA01} & Desconto percentual (\% sobre o subtotal) & \texttt{desconto = subtotal * (valor / 100)} \\
        \textbf{FA02} & Desconto em valor fixo (R\$) & \texttt{desconto = min(valor, subtotal)} \\
        \midrule
        \textbf{FE01} & Cupom inexistente ou inativo (\texttt{ativo=0}) & \texttt{CupomInexistenteError} / \texttt{CupomInvalidoError} \\
        \textbf{FE02} & Data atual ultrapassa o limite de validade & \texttt{CupomExpiradoError} \\
        \textbf{FE03} & Usos realizados atingiram o limite & \texttt{CupomEsgotadoError} \\
        \textbf{FE04} & Subtotal da compra menor que o mínimo exigido & \texttt{ValorMinimoNaoAtingidoError} \\
        \bottomrule
    \end{tabular}
    \end{table}
    \vspace{0.2cm}
    \begin{alertblock}{Conexão Direta: Especificação $\rightarrow$ Código}
    Cada fluxo de exceção (FE01 a FE04) mapeia exatamente para uma classe em \texttt{exceptions.py}!
    \end{alertblock}
\end{frame}

\section{Diagrama de Classes (UML 2)}

\begin{frame}{Modelagem Estrutural: Diagrama de Classes do Domínio}
\centering
\resizebox{0.92\textwidth}{!}{%
\begin{tikzpicture}[>=Latex, font=\sffamily]
    % Classe Cupom
    \node[draw=blue!80!black, fill=blue!5, thick, rounded corners=2pt, inner sep=3pt] (cupom) at (0,0) {
        \begin{tabular}{l}
        \textbf{\large Cupom} \\
        \hline
        \scriptsize - codigo: str \\
        \scriptsize - tipo\_desconto: TipoDesconto \\
        \scriptsize - valor\_desconto: float \\
        \scriptsize - valor\_minimo\_pedido: float \\
        \scriptsize - limite\_usos: int \\
        \scriptsize - usos\_realizados: int \\
        \scriptsize - validade\_ate: date \\
        \scriptsize - ativo: bool \\
        \hline
        \scriptsize + esta\_expirado(data\_ref): bool \\
        \scriptsize + atingiu\_limite\_usos(): bool \\
        \scriptsize + pode\_ser\_aplicado(subtotal): void \\
        \scriptsize + calcular\_desconto(subtotal): float \\
        \scriptsize + registrar\_uso(): void
        \end{tabular}
    };

    % Enum TipoDesconto
    \node[draw=purple!80!black, fill=purple!5, thick, rounded corners=2pt, inner sep=3pt, right=1.2cm of cupom, yshift=1.6cm] (enum) {
        \begin{tabular}{l}
        \textbf{$\ll$enumeration$\gg$} \\
        \textbf{TipoDesconto} \\
        \hline
        \scriptsize PERCENTUAL \\
        \scriptsize VALOR\_FIXO
        \end{tabular}
    };

    % Classe Pedido
    \node[draw=green!60!black, fill=green!5, thick, rounded corners=2pt, inner sep=3pt, right=1.2cm of cupom, yshift=-1.4cm] (pedido) {
        \begin{tabular}{l}
        \textbf{\large Pedido} \\
        \hline
        \scriptsize - id: int \\
        \scriptsize - cliente: Cliente \\
        \scriptsize - itens: List[ItemPedido] \\
        \scriptsize - cupom: Optional[Cupom] \\
        \scriptsize - desconto\_aplicado: float \\
        \hline
        \scriptsize + aplicar\_cupom(cupom): void \\
        \scriptsize + subtotal: float \\
        \scriptsize + total: float
        \end{tabular}
    };

    % Conexoes
    \draw[->, dashed, thick, draw=purple!80!black] (cupom) -- node[midway, above, font=\tiny\bfseries] {usa} (enum);
    \draw[thick, draw=black!70] (pedido) -- node[near start, above, font=\tiny\bfseries] {0..1} node[near end, above, font=\tiny\bfseries] {1..*} node[midway, below, font=\tiny\itshape] {possui} (cupom);
\end{tikzpicture}
}
\end{frame}

\begin{frame}{Visão Geral das Classes da Livraria BiblioTauá}
\centering
\resizebox{0.92\textwidth}{!}{%
\begin{tikzpicture}[>=Latex, font=\sffamily]
    % Classe Cliente
    \node[draw=gray!80!black, fill=gray!5, thick, rounded corners=2pt, inner sep=3pt] (cliente) at (0, 1.8) {
        \begin{tabular}{l}
        \textbf{Cliente} \\
        \hline
        \scriptsize - id: int \\
        \scriptsize - nome: str \\
        \scriptsize - email: str \\
        \scriptsize - matricula: str
        \end{tabular}
    };

    % Classe Pedido
    \node[draw=green!60!black, fill=green!5, thick, rounded corners=2pt, inner sep=3pt] (pedido) at (4.0, 1.8) {
        \begin{tabular}{l}
        \textbf{Pedido} \\
        \hline
        \scriptsize - id: int \\
        \scriptsize - status: str \\
        \scriptsize - total: float
        \end{tabular}
    };

    % Classe ItemPedido
    \node[draw=gray!80!black, fill=gray!5, thick, rounded corners=2pt, inner sep=3pt] (item) at (4.0, -1.0) {
        \begin{tabular}{l}
        \textbf{ItemPedido} \\
        \hline
        \scriptsize - quantidade: int \\
        \scriptsize - preco\_unitario: float \\
        \hline
        \scriptsize + subtotal: float
        \end{tabular}
    };

    % Classe Livro
    \node[draw=gray!80!black, fill=gray!5, thick, rounded corners=2pt, inner sep=3pt] (livro) at (0, -1.0) {
        \begin{tabular}{l}
        \textbf{Livro} \\
        \hline
        \scriptsize - id: int \\
        \scriptsize - titulo: str \\
        \scriptsize - estoque: int \\
        \scriptsize - preco: float
        \end{tabular}
    };

    % Relacionamentos
    \draw[thick] (cliente) -- node[midway, above, font=\tiny] {realiza 1..*} (pedido);
    \draw[thick] (pedido) -- node[midway, right, font=\tiny] {contém 1..*} (item);
    \draw[thick] (item) -- node[midway, above, font=\tiny] {refere 1} (livro);
\end{tikzpicture}
}
\end{frame}

\section{Arquitetura do Projeto BiblioTauá}

\begin{frame}[fragile]{Arquitetura do Projeto: Estrutura de Pastas e Módulos}
\begin{columns}[T]
    \begin{column}{0.53\textwidth}
        \begin{block}{\footnotesize Árvore de Arquivos do BiblioTauá}
\begin{lstlisting}[basicstyle=\fontsize{5.8pt}{6.8pt}\ttfamily]
bibliotaua/
|-- database.py          (SQLite v0.3.0)
|-- exceptions.py        (NOVO - Aula 24)
|-- models/
|   |-- cliente.py       (Base v0.3.0)
|   |-- livro.py         (Base v0.3.0)
|   |-- item_pedido.py   (Base v0.3.0)
|   |-- pedido.py        (Base v0.3.0)
|   `-- cupom.py         (NOVO - Aula 24)
`-- repositories/
    |-- livro_repo.py    (Base v0.3.0)
    `-- pedido_repo.py   (Base v0.3.0)
\end{lstlisting}
        \end{block}
    \end{column}
    \begin{column}{0.44\textwidth}
        \begin{block}{\footnotesize Papel de Cada Camada}
        \scriptsize
        \begin{itemize}
            \item \textbf{\texttt{models/}:} Regras de negócio e entidades OO puras.
            \item \textbf{\texttt{repositories/}:} Consultas SQL e persistência SQLite.
            \item \textbf{\texttt{database.py}:} Gerenciamento de conexão e DDL.
            \item \textbf{\texttt{exceptions.py}:} Exceções mapeadas da DUC.
        \end{itemize}
        \end{block}
    \end{column}
\end{columns}
\end{frame}

\begin{frame}{Por que o Repositório tem Mais Códigos que os Slides?}
    \begin{alertblock}{Simulação Real de Engenharia de Software Ágil}
    Em um time profissional, um novo ticket de Sprint nunca começa com o repositório em branco!
    \end{alertblock}
    \vspace{0.2cm}
    \begin{columns}[T]
        \begin{column}{0.48\textwidth}
            \begin{block}{\footnotesize 1. Base Pré-existente (v0.3.0)}
            \scriptsize
            \begin{itemize}
                \item Construída nas Sprints 01 a 03.
                \item Contém as entidades que já funcionam: \texttt{Cliente}, \texttt{Livro}, \texttt{ItemPedido}, \texttt{Pedido}.
                \item Evita erros de importação (\texttt{ModuleNotFoundError}) e viabiliza a execução do sistema completo.
            \end{itemize}
            \end{block}
        \end{column}
        \begin{column}{0.48\textwidth}
            \begin{block}{\footnotesize 2. O Nosso Foco Hoje (TASK-204)}
            \scriptsize
            \begin{itemize}
                \item Baixe o \textbf{\texttt{projeto\_aluno.zip}} com os arquivos base e marcações \texttt{TODO}.
                \item \textbf{\texttt{exceptions.py}:} Criar classes de erro da DUC.
                \item \textbf{\texttt{models/cupom.py}:} Criar métodos de regra.
                \item Teste criando o script \textbf{\texttt{main.py}} no projeto!
            \end{itemize}
            \end{block}
        \end{column}
    \end{columns}
\end{frame}

\section{Implementação: Exceções de Domínio}

\begin{frame}[fragile]{Mão na Massa: Hierarquia de Exceções (\texttt{exceptions.py})}
\begin{columns}[T]
    \begin{column}{0.48\textwidth}
\begin{lstlisting}[language=Python, basicstyle=\fontsize{5.2pt}{6.1pt}\ttfamily]
class BiblioTauaException(Exception):
    """Excecao base para regras de negocio."""
    pass

class CupomInvalidoError(BiblioTauaException):
    """Erro geral de cupom invalido."""
    pass

class CupomInexistenteError(CupomInvalidoError):
    def __init__(self, codigo: str):
        super().__init__(f"Cupom '{codigo}' nao encontrado no sistema.")
        self.codigo = codigo

class CupomInativoError(CupomInvalidoError):
    def __init__(self, codigo: str):
        super().__init__(f"O cupom '{codigo}' encontra-se desativado temporariamente.")
        self.codigo = codigo

class CupomExpiradoError(CupomInvalidoError):
    def __init__(self, codigo: str, validade_ate: str):
        super().__init__(f"O cupom '{codigo}' expirou em {validade_ate}.")
        self.codigo = codigo
        self.validade_ate = validade_ate
\end{lstlisting}
    \end{column}
    \begin{column}{0.50\textwidth}
\begin{lstlisting}[language=Python, basicstyle=\fontsize{5.2pt}{6.1pt}\ttfamily]
class CupomEsgotadoError(CupomInvalidoError):
    def __init__(self, codigo: str, limite: int):
        super().__init__(f"O cupom '{codigo}' atingiu a quantidade maxima de {limite} usos e esgotou.")
        self.codigo = codigo
        self.limite = limite

class ValorMinimoNaoAtingidoError(CupomInvalidoError):
    def __init__(self, codigo: str, valor_minimo: float, subtotal_atual: float):
        super().__init__(
            f"O cupom '{codigo}' exige pedido minimo de R$ {valor_minimo:.2f}. "
            f"Subtotal atual: R$ {subtotal_atual:.2f}."
        )
        self.codigo = codigo
        self.valor_minimo = valor_minimo
        self.subtotal_atual = subtotal_atual
\end{lstlisting}
    \end{column}
\end{columns}
\end{frame}

\section{Implementação: Modelo Cupom}

\begin{frame}[fragile]{Mão na Massa: Entidade \texttt{Cupom} (\texttt{models/cupom.py})}
\begin{lstlisting}[language=Python, basicstyle=\fontsize{4.8pt}{5.7pt}\ttfamily]
from dataclasses import dataclass
from datetime import datetime, date
from enum import Enum
from typing import Optional, Union
from bibliotaua.exceptions import *

class TipoDesconto(str, Enum):
    PERCENTUAL = "PERCENTUAL"
    VALOR_FIXO = "VALOR_FIXO"

@dataclass
class Cupom:
    codigo: str
    tipo_desconto: TipoDesconto
    valor_desconto: float
    valor_minimo_pedido: float = 0.0
    limite_usos: int = 100
    usos_realizados: int = 0
    validade_ate: Union[str, date] = "2026-12-31"
    ativo: bool = True
    id: Optional[int] = None

    def __post_init__(self):
        self.codigo = self.codigo.strip().upper()
        if not self.codigo:
            raise ValueError("O codigo do cupom nao pode ser vazio.")
        if self.valor_desconto <= 0:
            raise ValueError("O valor de desconto deve ser estritamente positivo.")
        if self.tipo_desconto == TipoDesconto.PERCENTUAL and self.valor_desconto > 100:
            raise ValueError("O desconto percentual nao pode ser superior a 100%.")
        if self.valor_minimo_pedido < 0:
            raise ValueError("O valor minimo do pedido nao pode ser negativo.")
        if self.limite_usos <= 0:
            raise ValueError("O limite de utilizacoes deve ser de no minimo 1.")
        if self.usos_realizados < 0:
            raise ValueError("A contagem de utilizacoes nao pode ser negativa.")
\end{lstlisting}
\end{frame}

\begin{frame}[fragile]{Mão na Massa: Métodos de Validação (\texttt{models/cupom.py})}
\begin{lstlisting}[language=Python, basicstyle=\fontsize{5.8pt}{6.8pt}\ttfamily]
    def obter_data_validade(self) -> date:
        if isinstance(self.validade_ate, str):
            return datetime.strptime(self.validade_ate, "%Y-%m-%d").date()
        return self.validade_ate

    def esta_expirado(self, data_referencia: Optional[date] = None) -> bool:
        hoje = data_referencia or date.today()
        return hoje > self.obter_data_validade()

    def atingiu_limite_usos(self) -> bool:
        return self.usos_realizados >= self.limite_usos

    def registrar_uso(self) -> None:
        if self.atingiu_limite_usos():
            raise CupomEsgotadoError(self.codigo, self.limite_usos)
        self.usos_realizados += 1
\end{lstlisting}
\end{frame}

\begin{frame}[fragile]{Mão na Massa: Aplicação e Desconto (\texttt{models/cupom.py})}
\begin{lstlisting}[language=Python, basicstyle=\fontsize{5.6pt}{6.6pt}\ttfamily]
    def pode_ser_aplicado_ao_pedido(self, subtotal: float, data_referencia: Optional[date] = None) -> None:
        if not self.ativo:
            raise CupomInativoError(self.codigo)
        if self.esta_expirado(data_referencia):
            validade_str = self.obter_data_validade().strftime("%d/%m/%Y")
            raise CupomExpiradoError(self.codigo, validade_str)
        if self.atingiu_limite_usos():
            raise CupomEsgotadoError(self.codigo, self.limite_usos)
        if subtotal < self.valor_minimo_pedido:
            raise ValorMinimoNaoAtingidoError(self.codigo, self.valor_minimo_pedido, subtotal)

    def calcular_desconto(self, subtotal: float, data_referencia: Optional[date] = None) -> float:
        self.pode_ser_aplicado_ao_pedido(subtotal, data_referencia)
        if self.tipo_desconto == TipoDesconto.PERCENTUAL:
            desconto = (subtotal * self.valor_desconto) / 100.0
        else:
            desconto = self.valor_desconto
        return round(min(desconto, subtotal), 2)
\end{lstlisting}
\end{frame}

\begin{frame}[fragile]{Testando a Implementação: Arquivo \texttt{main.py}}
\begin{columns}[T]
    \begin{column}{0.54\textwidth}
        \begin{block}{\footnotesize Criando o script de teste (\texttt{main.py})}
\begin{lstlisting}[language=Python, basicstyle=\fontsize{4.6pt}{5.5pt}\ttfamily]
from datetime import date, timedelta
from bibliotaua.models.cupom import Cupom, TipoDesconto
from bibliotaua.exceptions import ValorMinimoNaoAtingidoError

def main():
    print("--- Testando a Entidade Cupom (BiblioTaua) ---")

    # 1. Instanciando cupom de 10% (minimo R$ 50)
    cupom = Cupom(
        codigo="BEMVINDO10",
        tipo_desconto=TipoDesconto.PERCENTUAL,
        valor_desconto=10.0,
        valor_minimo_pedido=50.0,
        validade_ate=date.today() + timedelta(days=5)
    )
    print(f"Cupom: {cupom.codigo} ({cupom.valor_desconto}% OFF)")

    # 2. Teste valido (R$ 100.00)
    desconto = cupom.calcular_desconto(100.0)
    print(f"Subtotal R$ 100.00 -> Desconto: R$ {desconto:.2f}")

    # 3. Teste abaixo do minimo (R$ 30.00)
    try:
        cupom.calcular_desconto(30.0)
    except ValorMinimoNaoAtingidoError as erro:
        print(f"Subtotal R$ 30.00  -> Bloqueado: {erro}")

if __name__ == "__main__":
    main()
\end{lstlisting}
        \end{block}
    \end{column}
    \begin{column}{0.44\textwidth}
        \begin{block}{\footnotesize Execução no Terminal}
\begin{lstlisting}[basicstyle=\fontsize{4.6pt}{5.5pt}\ttfamily]
$ python3 main.py

--- Testando a Entidade Cupom (BiblioTaua) ---
Cupom: BEMVINDO10 (10.0% OFF)
Subtotal R$ 100.00 -> Desconto: R$ 10.00
Subtotal R$ 30.00  -> Bloqueado: O cupom
'BEMVINDO10' exige pedido minimo de R$ 50.00.
Subtotal atual: R$ 30.00.
\end{lstlisting}
        \end{block}
        \vspace{0.15cm}
        \begin{exampleblock}{\footnotesize Conclusão da Aula 24}
        \scriptsize
        A entidade \texttt{Cupom} e as regras de negócio foram validadas com sucesso pelo script \texttt{main.py}!
        \end{exampleblock}
    \end{column}
\end{columns}
\end{frame}

\section{Rito Ágil e Próximos Passos}

\begin{frame}{Rito Ágil: Registro da Daily 1}
    \begin{block}{Anotação para o Relatório da Sprint no \texttt{guia\_desenvolvedor.md}}
    \begin{itemize}
        \item \textbf{O que realizei hoje?}
        \begin{itemize}
            \item Interpretei o Diagrama de Casos de Uso e o Diagrama de Classes.
            \item Mapeei todos os fluxos da DUC UC04 em regras técnicas.
            \item Implementei a hierarquia de exceções em \texttt{exceptions.py}.
            \item Implementei a entidade de domínio \texttt{Cupom} e enum \texttt{TipoDesconto} com validações no \texttt{models/cupom.py}.
        \end{itemize}
        \item \textbf{O que farei no Encontro 2 (Aula 25)?}
        \begin{itemize}
            \item Mapeamento relacional (DER): criação da tabela \texttt{cupons} no SQLite.
            \item Implementação do padrão repositório em \texttt{CupomRepository}.
        \end{itemize}
        \item \textbf{Impedimentos?} Nenhum impedimento.
    \end{itemize}
    \end{block}
\end{frame}

\begin{frame}{Síntese do Encontro 1}
    \begin{itemize}
        \item \textbf{Objetivo Cumprido com Sucesso:} Saímos do alinhamento ágil, modelamos o sistema em UML (UC + Classes) e entregamos os primeiros arquivos de código em produção.
        \item \textbf{Para o Encontro 2 (Aula 25):}
        \begin{itemize}
            \item Persistência relacional com SQLite.
            \item Script DDL e constraints \texttt{CHECK}, \texttt{UNIQUE}.
            \item Repositório com prevenção contra SQL Injection.
        \end{itemize}
    \end{itemize}
    \vspace{0.4cm}
    \centering
    \textbf{Excelente início! O motor de cupons já começou a ganhar vida!}
\end{frame}

\end{document}
